% Part: proof-theory
% Chapter: normalization
% Section: introduction

\documentclass[../../../include/open-logic-section]{subfiles}

\begin{document}

\olfileid[pt]{pt}{nor}{seg}

\olsection{Segmentos e cortes}

!!^a{introduction} seguida por !!a{elimination} constitui um desvio em
!!a{proof}, e normalizar !!a{proof} é eliminar desvios desse tipo. No
entanto, \Elim{\lor} e \Elim{\lexists} tornam um pouco complicada até a
definição de ``desvio que queremos eliminar''. O problema é o seguinte.
Por causa dessas regras, não há garantia de que, em um desvio, a
!!{elimination} siga \emph{imediatamente} a !!{introduction}. Pode haver
uma \Elim{\lor} ou uma \Elim{\lexists} entre elas, por exemplo:
\[
\AxiomC{}
\DeduceC{$\lexists[x][!C(x)]$}
\AxiomC{$\Discharge{!A}{x}$}
\DeduceC{$!B$}
\DischargeRule{\Intro{\lif}}{x}
\UnaryInfC{$!A \lif !B$}
\DischargeRule{\Elim{\lexists}}{y}
\BinaryInfC{$!A \lif !B$}
\AxiomC{}
\DeduceC{$!A$}
\RightLabel{\Elim{\lif}}
\BinaryInfC{$!B$}
\DisplayProof
\]
Na verdade, qualquer número de \Elim{\lor} ou de
\Elim{\lexists} pode separar \Intro{\lif} de \Elim{\lif}, por exemplo:
\[
\AxiomC{}
\DeduceC{$\lexists[x][!C(x)]$}
\AxiomC{}
\DeduceC{$!D \lor !E$}
\AxiomC{$\Discharge{!A}{x}$}
\DeduceC{$!B$}
\DischargeRule{\Intro{\lif}}{x}
\UnaryInfC{$!A \lif !B$}
\AxiomC{$\Discharge{!A}{x}$}
\DeduceC{$!B$}
\DischargeRule{\Intro{\lif}}{x}
\UnaryInfC{$!A \lif !B$}
\RightLabel{\Elim{\lor}}
\TrinaryInfC{$!A \lif !B$}
\DischargeRule{\Elim{\lexists}}{y}
\BinaryInfC{$!A \lif !B$}
\AxiomC{}
\DeduceC{$!A$}
\RightLabel{\Elim{\lif}}
\BinaryInfC{$!B$}
\DisplayProof
\]
Esse exemplo mostra até que a mesma \Elim{\lif} pode estar envolvida em
vários desvios, pois sucede a mais de uma~\Intro{\lif}. Para levar em
conta todas essas possibilidades, definimos os desvios usando certos
tipos de sequências de ocorrências de !!{formula} em uma prova de
\Log{N1i}, chamadas \emph{segmentos}. Em nosso exemplo, trata-se de uma
sequência de ocorrências de~$!A \lif !B$ em que uma premissa menor de
$\Elim{\lor}$ ou de $\Elim{\lexists}$ sucede à conclusão correspondente.
No exemplo, há dois segmentos desse tipo, cada um com três ocorrências
de~$!A \lif !B$. Eis a definição oficial.

\begin{defn}
Um \emph{segmento} em uma !!{proof} de \Log{N1i}~$\delta$ é uma sequência
de ocorrências $!A_1$, \dots, $!A_n$ da mesma !!{formula}~$!A$ em
$\delta$, na qual
\begin{enumerate}
  \item $!A_1$ não é a conclusão de \Elim{\lor} nem de \Elim{\lexists};
  \item $!A_n$ não é uma premissa menor de \Elim{\lor} nem de
  \Elim{\lexists};
  \item Para $i = 1$, \dots,~$n-1$, $!A_i$ é uma premissa menor de
  \Elim{\lor} ou de \Elim{\lexists}, e $A_{i+1}$ é a conclusão dessa
  inferência.
\end{enumerate}
O \emph{comprimento} desse segmento é~$n$.
\end{defn}

Nem todo segmento é alvo de eliminação (um desvio). Os segmentos que são
recebem o nome de \emph{cortes}.

\begin{defn}
Um \emph{corte} é um segmento no qual $!A_n$ é a premissa maior de
!!a{elimination} e ocorre uma das seguintes condições: $n=1$ e
$!A_1 = !A_n$ é a conclusão de !!a{introduction} ou de~\FalseInt;
ou~$n>1$.
\end{defn}

Observe que um segmento de corte não precisa começar com a conclusão de
!!a{introduction}. Exige-se apenas que ele termine com a premissa maior
de !!a{elimination}. No entanto, um segmento de comprimento~$1$ precisa
ser a conclusão de uma regra de !!{introduction} para contar como corte.

\begin{defn}
O \emph{posto} de um corte é~$\depth{!A}$. O \emph{posto de
corte}~$\cutrank{\delta}$ de !!a{proof}~$\delta$ é o maior posto dos
cortes em~$\delta$. Um corte é \emph{máximo} em~$\delta$ se seu posto é
$\cutrank{\delta}$, isto é, se tem posto máximo. O \emph{comprimento de
corte} de~$\delta$ é a soma dos comprimentos de todos os seus cortes
máximos.
\end{defn}

\end{document}
